\section{Stateless weights, stateful store: memory transplant}
\label{sec:transplant}

The thesis of the construction is that the weights are stateless and the store is the
state: a single frozen model is turned into many different stateful agents purely by
which store it carries. We test this literally by transplanting stores (r86).

\paragraph{The store \emph{is} the state.} Take one frozen model and build two stores
from two disjoint fact sets (two ``conversations''). With the model held fixed, we vary
only which store is loaded and ask each store's query. \tabref{tab:swap} shows the
result is diagonal: the model answers a query at full recall when the matching store is
loaded and at \emph{chance} when the other store is loaded---swapping the store swaps
the memory, with no change to a single weight. A scrambled-store control is also at
chance. So the store carries the state in full; the weights carry none of it.

\begin{table}[t]
\centering
\begin{tabular}{lcc}
\toprule
loaded store $\rightarrow$ query & query set 1 & query set 2 \\
\midrule
store 1 & \recall{100} & $0.0\%$ \\
store 2 & $0.0\%$ & \recall{96.9} \\
\midrule
store 1, scrambled & $3.1\%$ (1 of 32 queries) & --- \\
\bottomrule
\end{tabular}
\caption{One frozen model, two stores. Recall is high on the diagonal (the store that
holds the queried facts) and at chance off it (the other store) and under scrambling.
The model's behaviour is set entirely by the loaded store: the store is the state.}
\label{tab:swap}
\end{table}

\paragraph{But the state is bound to its weights.} The store's keys are the writing
model's token embeddings, and retrieval is an exact match against the reading model's
query. So a store should be portable across a store-swap on the \emph{same} weights (as
above) but \emph{not} across different weights: a store written by model $B$ should not
read back on model $A$, even though the facts are physically present in it, because
$A$'s exact-match query is computed from $A$'s embeddings, not $B$'s. This is why the
persistence path (\S\ref{sec:method}) stamps each store with a model fingerprint and
refuses a foreign one by default. We force the transplant anyway to measure what
``refuse'' is protecting against: model $A$ reads model $B$'s store on $B$'s own facts
(and vice-versa), the store present and shape-compatible, the only change being the
weights. The result is unambiguous (r86, two independently-seeded $d{=}256$, $6$-layer
recall models). Each model reads its \emph{own} store at full recall---model $A$ at
\recall{98.8}, model $B$ at \recall{100}---but a store written by the \emph{other} model
reads back at \textbf{$0.0\%$ in both directions}, despite the facts being physically
present and the store shape-compatible. A store is a memory only for the model that wrote
it: its keys are that model's token embeddings, and a differently-seeded model's
exact-match query does not align with them. This is exactly what the fingerprint check
prevents, and it holds for the gated managed store as well as this plain-store pair.

\paragraph{Why this matters.} The two halves together give a concrete operational
picture. A store is a detachable, swappable unit of memory: one shared frozen model can
serve many users by hot-swapping their stores, each recovering its own memory exactly
(the diagonal of \tabref{tab:swap}), which is the multi-tenant version of ``memory as
a block.'' At the same time memory is not a universal artifact that can be moved between
models---it is encoded in the writing model's representation space, so porting it to a
different model requires re-encoding, not a file copy. Statelessness of the weights buys
the swappability; the exact-token key that makes retrieval work is exactly what binds a
store to its author.
